% ================================================================
% ==== FILE: content/cycles/cycles_k2.tex
% ================================================================

\subsubsection{Zyklen auf \texorpdfstring{$K_2$}{K_2}}
\label{sec:cycles-k2}

Ebene$K_2$entspricht dem Kontinuum der Konnektivität.
Sein Zustandsraum wird durch eine versickerungsartige Konfiguration dargestellt
\[
  \Omega(K_2) \cong \{0,1\}^{|E|},
\]
Wo$E$ist die Menge der Kanten und die Strukturvariable ist die
Berufswahrscheinlichkeit$p$.
Das grundlegende Phänomen auf dieser Ebene ist die Entstehung oder Zerstörung
der globalen Konnektivität.

Zyklen weiter$K_2$entstehen aufgrund von Schwankungen in der Clusterstruktur,
connectivity measures, and the effective time $\tau(K_2)$ associated
mit Korrelationsdynamik.

\subsubsection{Definition von Zyklen auf \texorpdfstring{$K_2$}{K_2}}

Ein Zyklus weiter$K_2$ist eine dynamische Flugbahn
\[
  C:[0,\tau]\to\Omega(K_2)
\]
so dass:
\[
  C(0)=C(\tau),
\]
und seine Entwicklung wird vom Operator gesteuert
\[
  \frac{dp}{dt}
  = F_2(p,\, A_2,\, P_2,\, J_2,\, \Theta_2),
\]
Wo:

\begin{itemize}
  \item $A_2$sind Konnektivitätsachsen (lokale „besetzte/unbesetzte“ Zustände),
  \item $P_2$sind Potenziale, die aus Clustergrößen, Korrelationslänge,
oder energieähnliche Parameter,
  \item $J_2$sind Ströme von Konnektivitätsänderungen,
  \item $\Theta_2$enthält die kritische Versickerungsschwelle$p_c$.
\end{itemize}

Somit ist ein Zyklus eine geschlossene Wiederholung von Konnektivitätskonfigurationen.

\subsubsection{Arten von Zyklen auf \texorpdfstring{$K_2$}{K_2}}

Es gibt drei Hauptklassen von Zyklen.

\paragraph{(1) Cluster–Rearrangement Cycles.}
Lokale Edge-Flip-Dynamik erzeugt oszillierendes Verhalten im Cluster
Struktur.
Diese Zyklen nähern sich nicht der Kritikalität und weisen eine begrenzte Korrelation auf
Länge.

\paragraph{(2) Subcritical Connectivity Cycles.}
Für$p<p_c$, kann das System in der Größe wiederkehrendes Verhalten erfahren
Verteilung endlicher Cluster:
\[
  C: \quad
  N(s,t) \to N(s,t+\tau),
\]
Wo$N(s)$ist die Anzahl der Clustergrößen$s$.
Diese Zyklen sind aufgrund des Fehlens eines Riesen strukturell stabil
Komponente.

\paragraph{(3) Critical–Near-Critical Cycles.}
For $p\approx p_c$, oscillations may bring the system repeatedly close
bis zur kritischen Perkolationsschwelle.
Diese Zyklen haben:

\begin{itemize}
  \item lange Korrelationslängen$\xi(p(t))$,
  \item hohe strukturelle Spannung$T_2(t)$,
  \item Empfindlichkeit gegenüber Schwellenwerten$\Theta_2$,
  \item Potenzial für phasenübergangsartige Exkursionen.
\end{itemize}

Dies ist die erste Ebene in der Hierarchie mit nicht trivialen Strukturen
Zyklen erscheinen.

\subsubsection{Strukturelle Spannung und Schwelleninteraktion}

Strukturelle Spannung an$K_2$ist gegeben durch:
\[
T_2(p)
  = f\bigl(\xi(p),\; \mathrm{Conn}(p),\;
       |p - p_c|\bigr),
\]
Wo$f$wächst, wenn sich das System der Kritikalität nähert.

Ein Zyklus kann nur dann periodisch sein, wenn:
\[
  T_2(t) < \Theta_{\mathrm{death}}
\]
für alle$t$.
Wenn die Schleife die Phasengrenze überschreitet
\[
p = p_c,
\]
Das System durchläuft einen Strukturwandel und der Kreislauf kann sich nicht schließen
reibungslos, es sei denn, die Dynamik kehrt zurück$p$über die gleiche Grenze.

Dies führt zunächst einmal zur Möglichkeit **kritischer Zyklen**
im OC von Schleifen, die mit Phasenübergängen interagieren.

\subsubsection{Metriken von Zyklen auf \texorpdfstring{$K_2$}{K_2}}

In Übereinstimmung mit Lauf ~ 9 werden die Zyklen wie folgt gemessen:

\[
L(C)
  = \int_0^\tau \|\dot{p}(t)\|\, dt,
\]
\[
  C_{\mathrm{eff}}
  = \frac{1}{L(C)} \int_0^\tau J_2(t)\cdot dA_2(t),
\]
\[
  S(C) = \min_{t\in[0,\tau]} d\bigl(\Omega(K_2),\partial\Omega(K_2)\bigr),
\]
Wo$d$hängt von der Nähe zur Kritikalität ab.

Für nahezu kritische Zyklen:
\[
  S(C) \to 0 \quad \text{as} \quad p\to p_c.
\]

Dies ist der Ursprung der „kritischen Verlangsamung“ und des „kritischen Zyklus“.
Degeneration.''

\subsubsection{Geburt der Zeit und nicht-trivialer Zyklen}

According to the formal definition of $\tau(K_2)$ (memory block~C.1–7),
Zeit weiter$K_2$wird durch die Wiederholungsperiode der Korrelation bestimmt
Strukturen:
\[
  \tau(K_2) \sim \frac{1}{\omega_{\mathrm{corr}}}.
\]

So geht es weiter$K_2$stehen in direktem Zusammenhang mit:

\begin{itemize}
  \item Entwicklung der Korrelationslänge$\xi(t)$,
  \item Zeitskalen der Clusterumordnung,
  \item Beschäftigungsflussdynamik$dp/dt$.
\end{itemize}

Dies stellt fest$K_2$als erste Ebene mit **dynamischem Temporal
Komplexität**.

\subsubsection{Rolle von \texorpdfstring{$K_2$}{K_2} Zyklen in der Hierarchie}

Zyklen weiter$K_2$sind der Vorläufer von:

\begin{itemize}
  \item chemische Reaktionszyklen auf$K_3$;
  \item Membran- und Stoffwechselzyklen auf$K_4$;
  \item Spike-Zyklen an$K_5$;
  \item kognitive Wiederholungsschleifen weiter$K_6$;
  \item institutionelle Zyklen auf$K_7$;
  \item Reproduktionszyklen auf Zivilisationsebene weiter$K_8$.
\end{itemize}

The essential mechanism introduced at $K_2$ is \emph{critical recurrence}
— die Möglichkeit, zyklisch mit Schwellenwerten zu interagieren.

Daher$K_2$ist das minimale Niveau, das die Ontologie von Continua zulässt
nichttriviale, schwellensensitive, strukturell komplexe Zyklen.
